% GENERATED FILE. Edit canonical lessons, then run npm run generate:books.
% canonical-source-hash: fnv1a64:522c8bbccd1c33e2
% canonical-lessons: TE-C365-vegu, TE-C365-droham, TE-C365-pratijna, TE-C365-adesincu, TE-C365-bahiskarincu, TE-R365-review-358-365

\chapter{Spies and Traitors}
\label{ch:te-spies-and-traitors}
\hlchaptermodality{365}

\begin{chapteropening}
\textbf{By the end of this chapter:} \emph{I can say a scout, a betrayal, a vow, to order and to banish.}

Complete the last lesson of chapter 365: I can say a scout, a betrayal, a vow, to order and to banish.
\end{chapteropening}

\section[\texorpdfstring{\te{వేగు} (vēgu)}{వేగు (vēgu)}]{\te{వేగు} (vēgu) — a scout; a spy}

\label{lesson:TE-C365-vegu}



\begin{quote}
Before the new one: say the Telugu for a sacrifice, then the Telugu for a brave man.
\end{quote}

\subsection*{You'll want to know: \te{వేగు}}
\textbf{\te{వేగు}} — \emph{vēgu} — \textquotedblleft{}a scout; a spy\textquotedblright{}.

\begin{quote}
\te{రాజు} \te{శత్రువు} \te{దగ్గరకు} \te{వేగులను} \te{పంపాడు}.
\end{quote}

\begin{quote}
\emph{rāju śatruvu daggaraku vēgulanu pampāḍu.}
\end{quote}

\begin{quote}
The king sent scouts to the enemy.
\end{quote}

Plural: \te{వేగులు}, \emph{vēgulu}.

The native word for a spy, beside the Sanskrit \te{గూఢచారి}.

\begin{cousinweb}
\textbf{\te{వేగు}} (\emph{vēgu}) is native Dravidian: DEDR 5533 lists Telugu \emph{vēgu}, \textquotedblleft{}espionage, a spy\textquotedblright{}, with Tamil \emph{vēy}, \textquotedblleft{}to spy out\textquotedblright{}.
\end{cousinweb}

\begin{grammarlens}[title={What you've built}]
One more word for telling what happened.
\end{grammarlens}

\subsection*{Your turn}
\begin{itemize}
\raggedright
  \item \emph{Say it:} \emph{vēgu}
  \item \emph{Say it:} \emph{vēgu}, once more
  \item \emph{Say it:} \emph{rāju śatruvu daggaraku vēgulanu pampāḍu.} — The king sent scouts to the enemy.
  \item \emph{Recall:} say the Telugu for a sacrifice, then the Telugu for a brave man, then say \emph{vēgu} again
\end{itemize}

\begin{tcolorbox}[breakable,colback=teal!4,colframe=teal!35!black,title={Before you move on}]
What does \te{వేగు} mean? (\textquotedblleft{}A scout; a spy\textquotedblright{}.) Say it once more.
\end{tcolorbox}
\section[\texorpdfstring{\te{ద్రోహం} (drōhaṁ)}{ద్రోహం (drōhaṁ)}]{\te{ద్రోహం} (drōhaṁ) — a betrayal; treachery}

\label{lesson:TE-C365-droham}



\begin{quote}
Before the new one: say the Telugu for a brave man, then the Telugu for a scout.
\end{quote}

\subsection*{You'll want to know: \te{ద్రోహం}}
\textbf{\te{ద్రోహం}} — \emph{drōhaṁ} — \textquotedblleft{}a betrayal; treachery\textquotedblright{}.

\begin{quote}
\te{సేనాధిపతి} \te{రాజుకు} \te{ద్రోహం} \te{చేశాడు}.
\end{quote}

\begin{quote}
\emph{sēnādhipati rājuku drōhaṁ cēśāḍu.}
\end{quote}

\begin{quote}
The commander betrayed the king.
\end{quote}

\te{ద్రోహం} \te{చేయు}, to betray; the one betrayed takes -\te{కు}.

\begin{cousinweb}
\textbf{\te{ద్రోహం}} (\emph{drōhaṁ}) is borrowed from Sanskrit \emph{droha}, \textquotedblleft{}injury, perfidy, treachery\textquotedblright{} (Monier-Williams).
\end{cousinweb}

\begin{grammarlens}[title={What you've built}]
One more word for telling what happened.
\end{grammarlens}

\subsection*{Your turn}
\begin{itemize}
\raggedright
  \item \emph{Say it:} \emph{drōhaṁ}
  \item \emph{Say it:} \emph{drōhaṁ}, once more
  \item \emph{Say it:} \emph{sēnādhipati rājuku drōhaṁ cēśāḍu.} — The commander betrayed the king.
  \item \emph{Recall:} say the Telugu for a brave man, then the Telugu for a scout, then say \emph{drōhaṁ} again
\end{itemize}

\begin{tcolorbox}[breakable,colback=teal!4,colframe=teal!35!black,title={Before you move on}]
What does \te{ద్రోహం} mean? (\textquotedblleft{}A betrayal; treachery\textquotedblright{}.) Say it once more.
\end{tcolorbox}
\section[\texorpdfstring{\te{ప్రతిజ్ఞ} (pratijña)}{ప్రతిజ్ఞ (pratijña)}]{\te{ప్రతిజ్ఞ} (pratijña) — a vow; a pledge}

\label{lesson:TE-C365-pratijna}



\begin{quote}
Before the new one: say the Telugu for a scout, then the Telugu for a betrayal.
\end{quote}

\subsection*{You'll want to know: \te{ప్రతిజ్ఞ}}
\textbf{\te{ప్రతిజ్ఞ}} — \emph{pratijña} — \textquotedblleft{}a vow; a pledge\textquotedblright{}.

\begin{quote}
\te{వీరుడు} \te{రాజ్యాన్ని} \te{కాపాడతానని} \te{ప్రతిజ్ఞ} \te{చేశాడు}.
\end{quote}

\begin{quote}
\emph{vīruḍu rājyānni kāpāḍatānani pratijña cēśāḍu.}
\end{quote}

\begin{quote}
The hero vowed to protect the kingdom.
\end{quote}

\te{ప్రతిజ్ఞ} \te{చేయు}, to take a vow.

\begin{cousinweb}
\textbf{\te{ప్రతిజ్ఞ}} (\emph{pratijña}) is borrowed from Sanskrit \emph{pratijñā}, \textquotedblleft{}a promise\textquotedblright{} (CDIAL 8563).
\end{cousinweb}

\begin{grammarlens}[title={What you've built}]
One more word for telling what happened.
\end{grammarlens}

\subsection*{Your turn}
\begin{itemize}
\raggedright
  \item \emph{Say it:} \emph{pratijña}
  \item \emph{Say it:} \emph{pratijña}, once more
  \item \emph{Say it:} \emph{vīruḍu rājyānni kāpāḍatānani pratijña cēśāḍu.} — The hero vowed to protect the kingdom.
  \item \emph{Recall:} say the Telugu for a scout, then the Telugu for a betrayal, then say \emph{pratijña} again
\end{itemize}

\begin{tcolorbox}[breakable,colback=teal!4,colframe=teal!35!black,title={Before you move on}]
What does \te{ప్రతిజ్ఞ} mean? (\textquotedblleft{}A vow; a pledge\textquotedblright{}.) Say it once more.
\end{tcolorbox}
\section[\texorpdfstring{\te{ఆదేశించు} (ādēśiñcu)}{ఆదేశించు (ādēśiñcu)}]{\te{ఆదేశించు} (ādēśiñcu) — to order; to command}

\label{lesson:TE-C365-adesincu}



\begin{quote}
Before the new one: say the Telugu for a betrayal, then the Telugu for a vow.
\end{quote}

\subsection*{You'll want to know: \te{ఆదేశించు}}
\textbf{\te{ఆదేశించు}} — \emph{ādēśiñcu} — \textquotedblleft{}to order; to command\textquotedblright{}.

\begin{quote}
\te{రాజు} \te{ద్రోహిని} \te{పట్టుకోమని} \te{ఆదేశించాడు}.
\end{quote}

\begin{quote}
\emph{rāju drōhini paṭṭukōmani ādēśiñcāḍu.}
\end{quote}

\begin{quote}
The king ordered the traitor to be caught.
\end{quote}

Past (I …): \te{ఆదేశించాను}, \emph{ādēśiñcānu}.

Its noun is \te{ఆదేశం}, an order.

\begin{cousinweb}
\textbf{\te{ఆదేశించు}} (\emph{ādēśiñcu}) is built on Sanskrit \emph{ādeśa}, \textquotedblleft{}advice, instruction\textquotedblright{} (Monier-Williams), with the Telugu verb ending \emph{-iñcu}.
\end{cousinweb}

\begin{grammarlens}[title={What you've built}]
One more word for telling what happened.
\end{grammarlens}

\subsection*{Your turn}
\begin{itemize}
\raggedright
  \item \emph{Say it:} \emph{ādēśiñcu}
  \item \emph{Say it:} \emph{ādēśiñcu}, once more
  \item \emph{Say it:} \emph{rāju drōhini paṭṭukōmani ādēśiñcāḍu.} — The king ordered the traitor to be caught.
  \item \emph{Recall:} say the Telugu for a betrayal, then the Telugu for a vow, then say \emph{ādēśiñcu} again
\end{itemize}

\begin{tcolorbox}[breakable,colback=teal!4,colframe=teal!35!black,title={Before you move on}]
What does \te{ఆదేశించు} mean? (\textquotedblleft{}To order; to command\textquotedblright{}.) Say it once more.
\end{tcolorbox}
\section[\texorpdfstring{\te{బహిష్కరించు} (bahiṣkariñcu)}{బహిష్కరించు (bahiṣkariñcu)}]{\te{బహిష్కరించు} (bahiṣkariñcu) — to banish; to boycott}

\label{lesson:TE-C365-bahiskarincu}



\begin{quote}
Before the new one: say the Telugu for a vow, then the Telugu for to order.
\end{quote}

\subsection*{You'll want to know: \te{బహిష్కరించు}}
\textbf{\te{బహిష్కరించు}} — \emph{bahiṣkariñcu} — \textquotedblleft{}to banish; to boycott\textquotedblright{}.

\begin{quote}
\te{రాజు} \te{ఆ} \te{ద్రోహిని} \te{రాజ్యం} \te{నుంచి} \te{బహిష్కరించాడు}.
\end{quote}

\begin{quote}
\emph{rāju ā drōhini rājyaṁ nuñci bahiṣkariñcāḍu.}
\end{quote}

\begin{quote}
The king banished that traitor from the kingdom.
\end{quote}

Past (I …): \te{బహిష్కరించాను}, \emph{bahiṣkariñcānu}.

Also to boycott: \te{విద్యార్థులు} \te{తరగతులు} \te{బహిష్కరించారు}.

\begin{cousinweb}
\textbf{\te{బహిష్కరించు}} (\emph{bahiṣkariñcu}) is built on Sanskrit \emph{bahiṣkāra}, \textquotedblleft{}expulsion\textquotedblright{} (Monier-Williams), with the Telugu verb ending \emph{-iñcu}.
\end{cousinweb}

\begin{grammarlens}[title={What you've built}]
One more word for telling what happened.
\end{grammarlens}

\subsection*{Your turn}
\begin{itemize}
\raggedright
  \item \emph{Say it:} \emph{bahiṣkariñcu}
  \item \emph{Say it:} \emph{bahiṣkariñcu}, once more
  \item \emph{Say it:} \emph{rāju ā drōhini rājyaṁ nuñci bahiṣkariñcāḍu.} — The king banished that traitor from the kingdom.
  \item \emph{Recall:} say the Telugu for a vow, then the Telugu for to order, then say \emph{bahiṣkariñcu} again
\end{itemize}

\begin{tcolorbox}[breakable,colback=teal!4,colframe=teal!35!black,title={Before you move on}]
What does \te{బహిష్కరించు} mean? (\textquotedblleft{}To banish; to boycott\textquotedblright{}.) Say it once more.
\end{tcolorbox}
\section[(dialogue)]{Words from chapters 358–365 — a review}

\label{lesson:TE-R365-review-358-365}



\begin{quote}
Say the Telugu for to live and to settle down.
\end{quote}

\subsection*{Your turn}
Say these aloud:

\begin{itemize}
\raggedright
  \item to live, to settle down, to take shelter, to dream of, luckily
  \item an empire, an emperor, a prince, a princess, to found
  \item a sword, armour, a weapon, a bow, a gun
  \item a battlefield, a troop, to besiege, to conquer, a defeat
  \item to battle, to drive away, to bow one's head, a truce, an envoy
  \item a rebellion, to rise up, to suppress, to incite, independence
  \item a slave, liberation, a struggle, a sacrifice, a brave man
  \item a scout, a betrayal, a vow, to order, to banish
\end{itemize}

\begin{tcolorbox}[breakable,colback=teal!4,colframe=teal!35!black,title={Before you move on}]
To live? (\textbf{\te{జీవించు}}, \emph{jīviñcu}.) To settle down? (\textbf{\te{స్థిరపడు}}, \emph{sthirapaḍu}.) To take shelter? (\textbf{\te{తలదాచుకో}}, \emph{taladācukō}.) To dream of? (\textbf{\te{కలగను}}, \emph{kalaganu}.) Luckily? (\textbf{\te{అదృష్టవశాత్తు}}, \emph{adṛṣṭavaśāttu}.) An empire? (\textbf{\te{సామ్రాజ్యం}}, \emph{sāmrājyaṁ}.) An emperor? (\textbf{\te{చక్రవర్తి}}, \emph{cakravarti}.) A prince? (\textbf{\te{రాకుమారుడు}}, \emph{rākumāruḍu}.) A princess? (\textbf{\te{రాకుమారి}}, \emph{rākumāri}.) To found? (\textbf{\te{స్థాపించు}}, \emph{sthāpiñcu}.) A sword? (\textbf{\te{ఖడ్గం}}, \emph{khaḍgaṁ}.) Armour? (\textbf{\te{కవచం}}, \emph{kavacaṁ}.) A weapon? (\textbf{\te{ఆయుధం}}, \emph{āyudhaṁ}.) A bow? (\textbf{\te{విల్లు}}, \emph{villu}.) A gun? (\textbf{\te{తుపాకి}}, \emph{tupāki}.) A battlefield? (\textbf{\te{యుద్ధభూమి}}, \emph{yuddhabhūmi}.) A troop? (\textbf{\te{దళం}}, \emph{daḷaṁ}.) To besiege? (\textbf{\te{ముట్టడించు}}, \emph{muṭṭaḍiñcu}.) To conquer? (\textbf{\te{జయించు}}, \emph{jayiñcu}.) A defeat? (\textbf{\te{పరాజయం}}, \emph{parājayaṁ}.) To battle? (\textbf{\te{పోరాడు}}, \emph{pōrāḍu}.) To drive away? (\textbf{\te{తరిమికొట్టు}}, \emph{tarimikoṭṭu}.) To bow one's head? (\textbf{\te{తలవంచు}}, \emph{talavañcu}.) A truce? (\textbf{\te{సంధి}}, \emph{sandhi}.) An envoy? (\textbf{\te{దూత}}, \emph{dūta}.) A rebellion? (\textbf{\te{తిరుగుబాటు}}, \emph{tirugubāṭu}.) To rise up? (\textbf{\te{తిరగబడు}}, \emph{tiragabaḍu}.) To suppress? (\textbf{\te{అణచివేయు}}, \emph{aṇacivēyu}.) To incite? (\textbf{\te{ఉసిగొల్పు}}, \emph{usigolpu}.) Independence? (\textbf{\te{స్వాతంత్ర్యం}}, \emph{svātantryaṁ}.) A slave? (\textbf{\te{బానిస}}, \emph{bānisa}.) Liberation? (\textbf{\te{విముక్తి}}, \emph{vimukti}.) A struggle? (\textbf{\te{పోరాటం}}, \emph{pōrāṭaṁ}.) A sacrifice? (\textbf{\te{త్యాగం}}, \emph{tyāgaṁ}.) A brave man? (\textbf{\te{వీరుడు}}, \emph{vīruḍu}.) A scout? (\textbf{\te{వేగు}}, \emph{vēgu}.) A betrayal? (\textbf{\te{ద్రోహం}}, \emph{drōhaṁ}.) A vow? (\textbf{\te{ప్రతిజ్ఞ}}, \emph{pratijña}.) To order? (\textbf{\te{ఆదేశించు}}, \emph{ādēśiñcu}.) To banish? (\textbf{\te{బహిష్కరించు}}, \emph{bahiṣkariñcu}.)
\end{tcolorbox}
